\section{Bernoulli compensation and negative spectral orders}
\label{sec:compensation}

\subsection{A hierarchy with an exact convergence domain}
The basic kernel $h$ is bounded but does not vanish at zero. Its Mellin
transform therefore does not directly reach the Gamma logarithm at spectral
order zero. We now subtract exactly the required Taylor terms.
Let $B_{2k}$ be the Bernoulli numbers and put
\begin{align}
 P_0(x)&=x^{-1},\label{eq:P0}\\
 P_d(x)&=x^{-1}+\frac12+
       \sum_{k=1}^{d-1}\frac{B_{2k}}{(2k)!}x^{2k-1},\qquad d\ge1,
                 \label{eq:Pd}\\
 h_d(x)&=Q(x)-P_d(x),\qquad
 \nu_d=\begin{cases}0,&d=0,\\2d-1,&d\ge1.\end{cases}
                 \label{eq:hd}
\end{align}
An empty sum is zero. Thus $h_0=h$, $h_1=h-1/2$, and
$h_2=h-1/2-x/12$. The convergent local Bernoulli expansion gives
\begin{equation}\label{eq:hd-valuation}
 h_0(x)=\tfrac12+O(x),\qquad
 h_d(x)=\frac{B_{2d}}{(2d)!}x^{2d-1}+O(x^{2d+1})\quad(d\ge1).
\end{equation}
Write $P_d(x)=\sum_{j\in E_d}\alpha_{d,j}x^j$, where
$E_0=\{-1\}$ and
$E_d=\{-1,0,1,3,\ldots,2d-3\}$ for $d\ge1$, with the evident empty
positive portion for $d=1$.

\begin{definition}\label{def:Fd}
The Gamma-normalized compensated Hurwitz family is
\begin{equation}\label{eq:Fd-laplace}
 \mathcal F_d(s,z)=\int_0^\infty x^{s-1}e^{-zx}h_d(x)\dd x,
 \qquad\Rea s>-\nu_d,\quad\Rea z>0.
\end{equation}
Equivalently, on this domain it is the combined continuation of
\begin{equation}\label{eq:Fd-closed}
 \mathcal F_d(s,z)=\Gamma(s)\zeta(s,z)
       -\sum_{j\in E_d}\alpha_{d,j}\Gamma(s+j)z^{-s-j}.
\end{equation}
We also set $\mathcal R_d(s,z)=\mathcal F_d(s,z)/\Gamma(s)$ by
continuation. In particular,
\begin{align}
 \mathcal R_0(s,z)&=Z_s(z),\notag\\
 \mathcal R_d(s,z)&=Z_s(z)-\tfrac12z^{-s}
  -\sum_{k=1}^{d-1}\frac{B_{2k}}{(2k)!}(s)_{2k-1}z^{1-s-2k}
          \quad(d\ge1).\label{eq:Rd}
\end{align}
\end{definition}

\begin{proposition}[Holomorphy, zeros, and the position ladder]
\label{thm:Fd-domain}
The family $\mathcal F_d$ is jointly holomorphic on the domain in
\eqref{eq:Fd-laplace}. It is a holomorphic $L^2$ vector on the vertical
line $z=a+\ii y$ whenever
\begin{equation}\label{eq:Fd-L2}
 \Rea s>\frac12-\nu_d,\qquad\Rea a>0.
\end{equation}
It satisfies
\begin{equation}\label{eq:Fd-ladder}
 \partial_z\mathcal F_d(s,z)=-\mathcal F_d(s+1,z).
\end{equation}
For every integer $n$ with $0\le n<\nu_d$,
\begin{equation}\label{eq:Rd-zero}
 \mathcal R_d(-n,z)=0,\qquad
 \mathcal F_d(-n,z)=\frac{(-1)^n}{n!}
                    \left.\partial_s\mathcal R_d(s,z)\right|_{s=-n}.
\end{equation}
\end{proposition}
\begin{proof}
Near zero the kernel in \eqref{eq:Fd-laplace} is
$O(x^{\Rea s+\nu_d-1})$, and at infinity the exponential dominates
all polynomial terms in $P_d$. These bounds, with extra powers of
$\log x$ on compact spectral subsets, prove holomorphy. Squaring the
near-zero kernel gives exactly the sufficient $L^2$ condition
\eqref{eq:Fd-L2}. Plancherel then transfers the vector-valued holomorphy
to the vertical line. Multiplication by $-x$ proves the position ladder.

Formula~\eqref{eq:Fd-closed} follows by evaluating the separate terms
first for sufficiently large $\Rea s$, then continuing.
The function in \eqref{eq:Rd} is entire in $s$, after the original
$s=1$ pole cancellation. Since $1/\Gamma(s)$ has a simple zero at
$s=-n$ while $\mathcal F_d$ is holomorphic there, $\mathcal R_d(-n,z)=0$.
Finally,
$\Gamma(-n+\epsilon)=(-1)^n/(n!\epsilon)+O(1)$ gives the second formula.
\end{proof}

The kernel $L^2$ threshold is exact for an individual transform in its
Laplace domain, because the leading coefficients in
\eqref{eq:hd-valuation} are nonzero. We do not claim that the rectangular
pair of such thresholds is the maximal possible region for every product.
The theorem deliberately states a uniform domain in which all orders of
parameter differentiation are justified.

\subsection{Finite closure at every compensation depth}
Define
\begin{equation}\label{eq:J-integral}
 \mathcal J_{d,e}(w,A)=\int_0^\infty
          x^{w-1}e^{-Ax}h_d(x)h_e(x)\dd x.
\end{equation}
It is holomorphic for
$\Rea w>-\nu_d-\nu_e$, $\Rea A>0$.

\begin{theorem}[Arbitrary-depth product closure]\label{thm:depth-closure}
For $\Rea s>1/2-\nu_d$, $\Rea t>1/2-\nu_e$, and
$\Rea a,\Rea b>0$,
\begin{equation}\label{eq:Fd-pair}
 \frac1{2\pi}\int_{\R}\mathcal F_d(s,a+\ii y)
                         \mathcal F_e(t,b-\ii y)\dd y
       =\mathcal J_{d,e}(s+t-1,a+b).
\end{equation}
The right-hand side has the finite expression
\begin{align}
 \mathcal J_{d,e}(w,A)={}&
 \Gamma(w)\{\zeta(w-1,A)+(1-A)\zeta(w,A)\}\notag\\
 &-\sum_{j\in E_d}\alpha_{d,j}\Gamma(w+j)\zeta(w+j,A)\notag\\
 &-\sum_{k\in E_e}\alpha_{e,k}\Gamma(w+k)\zeta(w+k,A)\notag\\
 &+\sum_{j\in E_d}\sum_{k\in E_e}
   \alpha_{d,j}\alpha_{e,k}\Gamma(w+j+k)A^{-w-j-k}.
                  \label{eq:J-finite}
\end{align}
Every apparent pole of this expression in
$\Rea w>-\nu_d-\nu_e$ is removable. The identity and all its spectral
and position derivatives are identities of ordinary convergent pairings.
\end{theorem}
\begin{proof}
Equation~\eqref{eq:Fd-pair} is \cref{thm:plancherel} applied to
\eqref{eq:Fd-laplace}. For $\Rea w>2$, expand
\[
 (Q-P_d)(Q-P_e)=Q^2-P_dQ-P_eQ+P_dP_e.
\]
Every individual Mellin integral then converges. The $Q^2$ term was
evaluated earlier; the monomial--$Q$ terms yield
$\Gamma(w+j)\zeta(w+j,A)$, and the monomial products yield the last
line of \eqref{eq:J-finite}. The original combined integrand vanishes
to order $\nu_d+\nu_e$ at zero, proving its holomorphy on the larger
half-plane. Analytic continuation therefore cancels \emph{all} of the
apparent poles there, not merely those checked at finitely many indices.
The vector-valued argument proves the derivative assertion.
\end{proof}

For $d=e=0$, \eqref{eq:J-finite} is exactly
$\mathcal J_{0,0}(w,A)=\Gamma(w)K(w,A)$. Higher depths do not replace the
special-function alphabet: they enlarge its finite range of spectral
arguments while improving the admissible left half-plane. This is the
mechanism that permits antiderivatives to be paired without a divergent
vertical integral.

\section{Stirling remainders and normalized antiderivatives}
\label{sec:primitives}

\subsection{The first negative-order pairing}
At depth one, \eqref{eq:Rd-zero} and the standard zeta values at zero give
\begin{equation}\label{eq:Omega-F}
 \mathcal F_1(0,z)=\Omega(z)
    =\int_0^\infty e^{-zx}\frac{h_1(x)}x\dd x.
\end{equation}
For completeness, differentiate
$\mathcal R_1(s,z)=\zeta(s,z)-z^{1-s}/(s-1)-\frac12z^{-s}$ at $s=0$.
The result is
$\zeta'(0,z)-z\log z+z+\frac12\log z$, which is
\eqref{eq:Omega-intro}. Thus \eqref{eq:Omega-F} also follows directly
from the compensated Hurwitz family, independently of quoting Binet's
formula; it agrees with the classical representation in
\cite{dlmf-gamma}.

\begin{theorem}[Exact Stirling-remainder pairing]\label{thm:Omega-pair}
Let $\Rea a,\Rea b>0$ and $A=a+b$. Then
\begin{align}
 \frac1{2\pi}\int_{\R}\Omega(a+\ii y)\Omega(b-\ii y)\dd y
   ={}&A\zeta'(-1,A)-2\zeta'(-2,A)\notag\\
     &+\left(\frac{A^3}{6}-\frac{A^2}{2}+\frac A4\right)\log A
           +\frac{A^3}{36}+\frac A{12}.
                  \label{eq:Omega-pair}
\end{align}
\end{theorem}
\begin{proof}
The domain condition \eqref{eq:Fd-L2} holds at $d=1$, $s=0$.
Hence the left-hand side equals $\mathcal J_{1,1}(-1,A)$.
Since $h_1^2=h^2-h+1/4$, write
$\mathcal J_{1,1}(w,A)=\Gamma(w)\mathcal B(w,A)$, where
\begin{align}
 \mathcal B(w,A)={}&\frac{w-3}{w-1}\zeta(w-1,A)-A\zeta(w,A)
       +\frac{A^{2-w}}{(w-1)(w-2)}\notag\\
       &+\frac{A^{1-w}}{w-1}+\frac14A^{-w}.
                      \label{eq:B-Stirling}
\end{align}
At $w=-1$,
\[
 \mathcal B(-1,A)=2\zeta(-2,A)-A\zeta(-1,A)
                       +\frac{A^3}{6}-\frac{A^2}{2}+\frac A4=0.
\]
This is the Bernoulli-polynomial identity
$\zeta(-n,A)=-B_{n+1}(A)/(n+1)$ at $n=1,2$.
Because $\Res_{w=-1}\Gamma(w)=-1$,
\[
 \mathcal J_{1,1}(-1,A)=-\partial_w\mathcal B(-1,A).
\]
Differentiation before simplification gives
\begin{align*}
 -\partial_w\mathcal B(-1,A)={}&
 A\zeta'(-1,A)-2\zeta'(-2,A)-\tfrac12\zeta(-2,A)\\
 &+A^3\left(\frac{\log A}{6}-\frac5{36}\right)
       -\frac{A^2\log A}{2}+\frac{A^2}{4}+\frac{A\log A}{4}.
\end{align*}
Finally, $\zeta(-2,A)=-A^3/3+A^2/2-A/6$ yields
\eqref{eq:Omega-pair}.
\end{proof}

At $A=1$, the functional equation of the Riemann zeta function gives
$\zeta'(-2)=-\zeta(3)/(4\pi^2)$, proving
\eqref{eq:headline-stirling}. The sign in front of $\partial_w\mathcal B$
is forced by the negative residue of $\Gamma$ at $-1$.
Discarding the zero of $\mathcal B$ before differentiating would incorrectly
turn a nonzero norm into zero.

\subsection{A canonical primitive at every admissible depth}
Define the depth-$d$ Stirling remainder, for $d\ge1$, by
\begin{equation}\label{eq:Omega-depth}
 \Omega_d(z)=\Omega(z)-\sum_{k=1}^{d-1}
          \frac{B_{2k}}{2k(2k-1)}z^{1-2k}.
\end{equation}
The Gamma recurrence in \eqref{eq:Fd-closed} gives
$\mathcal F_d(0,z)=\Omega_d(z)$.

\begin{theorem}[Normalized primitive ladder and exact pairings]
\label{thm:primitive-ladder}
For $d\ge1$ and $0\le n<\nu_d=2d-1$, let
\begin{equation}\label{eq:primitive-definition}
 \mathcal P_{d,n}(z)=(-1)^n\mathcal F_d(-n,z)
      =\frac1{n!}\left.\partial_s\mathcal R_d(s,z)\right|_{s=-n}.
\end{equation}
Then
\begin{equation}\label{eq:primitive-ladder}
 \mathcal P_{d,0}=\Omega_d,\qquad
 \partial_z\mathcal P_{d,n}=\mathcal P_{d,n-1}\quad(n\ge1).
\end{equation}
Each $\mathcal P_{d,n}(x)$ tends to zero as real $x\to+\infty$.
This normalization uniquely specifies it among the holomorphic
$n$-fold antiderivatives of $\Omega_d$.
For $d,e\ge1$, $0\le n<\nu_d$, $0\le m<\nu_e$,
\begin{equation}\label{eq:primitive-pair}
 \frac1{2\pi}\int_{\R}\mathcal P_{d,n}(a+\ii y)
                       \mathcal P_{e,m}(b-\ii y)\dd y
       =(-1)^{n+m}\mathcal J_{d,e}(-n-m-1,a+b).
\end{equation}
The right-hand side is the removable value of the finite expression
\eqref{eq:J-finite}; no new integral is left unevaluated.
\end{theorem}
\begin{proof}
The definition and ladder follow from \cref{thm:Fd-domain}. The Laplace
representation is
\[
 \mathcal P_{d,n}(z)=(-1)^n\int_0^\infty x^{-n-1}e^{-zx}h_d(x)\dd x.
\]
Its near-zero power is $x^{\nu_d-n-1}$, with $\nu_d-n>0$.
Rescaling $x$ by a positive real $z$ and using an integrable split of the
kernel proves decay $O(z^{n-\nu_d})$ as $z\to+\infty$.
The difference of two $n$-fold antiderivatives is a polynomial of degree
less than $n$, and such a polynomial cannot tend to zero unless it is zero.
For $n=0$ the uniqueness statement simply identifies the function itself.

The integer restriction $n<\nu_d$ implies
$-n>1/2-\nu_d$; similarly for $m$. Thus \cref{thm:depth-closure} applies
and proves \eqref{eq:primitive-pair}. Moreover,
$-n-m-1>-\nu_d-\nu_e$, placing the spectral value strictly inside the
holomorphic domain of $\mathcal J_{d,e}$.
\end{proof}

The primitive formula is explicit in ordinary Hurwitz order derivatives:
\eqref{eq:Rd} involves only $\zeta(s,z)$, elementary powers, and a finite
Pochhammer polynomial. At $s=-n$ its first derivative therefore involves
$\zeta'(-n,z)$ and elementary terms. Pairing primitives at arbitrary
depths remains finite by \eqref{eq:J-finite}. At its integer resonances,
first derivatives of Hurwitz zeta and digamma values suffice; the next
section gives a stable finite-part rule. Higher \emph{spectral} jets may
require correspondingly higher zeta derivatives and Stieltjes constants,
but always only finitely many.
